% Partie : logique du premier ordre
% Chapitre : complétude
% Section : Löwenheim--Skolem descendant

\documentclass[../../../include/open-logic-section]{subfiles}

\begin{document}

\olfileid[fr]{fol}{com}{dls}
\olsection{Le théorème de Löwenheim--Skolem}

Dans le langage au plus dénombrable fixé ici, le théorème de
Löwenheim--Skolem affirme notamment que toute théorie ayant un
modèle infini possède aussi un modèle !!{enumerable}.
Il en découle que tout ensemble d'!!{sentence}s qui admet un
modèle !!{nonenumerable} admet aussi un modèle !!{enumerable}.
Un ensemble d'!!{sentence}s ne peut donc imposer que son
!!{domain} soit !!{nonenumerable}.

\begin{thm}
\ollabel{thm:downward-ls}
Si $\Gamma$ est cohérent, il possède un modèle !!{enumerable},
c'est-à-dire qu'il est satisfaisable dans !!a{structure} dont le
domaine est fini ou !!{denumerable}.
\end{thm}

\begin{proof}
Si $\Gamma$ est cohérent, la !!{structure}~$\Struct M$ construite
dans la démonstration du théorème de complétude a un domaine dont la
taille ne dépasse pas celle de l'ensemble des termes clos du
langage enrichi~$\Lang L'$. Ce domaine est formé soit des termes
clos, soit de leurs classes d'équivalence. Comme $\Lang L'$ est
au plus dénombrable, il en va de même de l'ensemble de ses termes
clos et du domaine de~$\Struct M$. Le réduit au langage initial
a le même domaine.
\end{proof}

\begin{thm}
\ollabel{noidentity-ls}
Si $\Gamma$ est un ensemble cohérent d'!!{sentence}s dans le
langage du premier ordre sans identité, il possède un modèle
!!{denumerable}, c'est-à-dire qu'il est satisfaisable dans !!a{structure}
dont le domaine est infini et !!{enumerable}.
\end{thm}

\begin{proof}
Si $\Gamma$ est cohérent et qu'aucun de ses énoncés ne contient
l'identité, la !!{structure}~$\Struct M$ construite dans la
démonstration de la complétude a pour domaine exactement l'ensemble
des \emph{termes clos} du langage enrichi~$\Lang L'$.
Cet ensemble est infini dénombrable, puisque le langage est au
plus dénombrable et que les constantes ajoutées par l'expansion
de Henkin sont des termes clos distincts. La structure~$\Struct M$
est donc !!{denumerable}.
\end{proof}

\begin{ex}[Paradoxe de Skolem]
La théorie des ensembles de Zermelo--Fraenkel avec axiome du choix,
$\Log{ZFC}$, est un
cadre assez puissant pour exprimer pratiquement tous les énoncés
mathématiques, y compris ceux qui portent sur la taille des
ensembles. Ainsi, $\Log{ZFC}$ démontre que l'ensemble~$\Real$ des
nombres réels est !!{nonenumerable}, et aussi le théorème de Cantor
selon lequel l'ensemble des parties de tout ensemble est plus grand
que cet ensemble lui-même. Si $\Log{ZFC}$ est cohérent, tous ses
modèles sont infinis et contiennent des !!{element}s que la théorie
décrit comme !!{nonenumerable}, par exemple celui qui représente
l'ensemble des parties des entiers naturels. Pourtant, d'après le
théorème de Löwenheim--Skolem, $\Log{ZFC}$ possède aussi des
modèles !!{enumerable} : vus de l'extérieur, ces modèles sont au
plus dénombrables, bien qu'ils contiennent des ensembles «
!!{nonenumerable}s » au sens interne du modèle. De l'extérieur, on
peut énumérer les objets que le modèle considère comme appartenant à
son ensemble des parties de~$\omega$. La non-dénombrabilité interne
exclut une bijection représentée dans le modèle entre cet ensemble
et ses entiers.
\end{ex}

\end{document}
